\documentclass[11pt]{article}
\usepackage[a4paper,margin=1in]{geometry}
\usepackage{amsmath,amsthm,amssymb}
\usepackage{microtype}
\usepackage[backend=biber,style=alphabetic]{biblatex}
\addbibresource{refs.bib}

\title{Related Work \\ \large{[draft, standalone for review]}}
\author{}
\date{}

\begin{document}
\maketitle

\noindent
This document drafts the Related Work section for the paper
\emph{Step-Indexed Hoare Logic as Ramified Forcing}. It is organized
into seven literature threads: (i) step-indexing in programming
language semantics; (ii) forcing in set theory, ramified and
unramified; (iii) topos-theoretic models of step-indexing; (iv)
forcing translations in type theory; (v) the modal logic of
provability (\textsf{GL}) and well-foundedness; (vi) constructive
modal logic for guarded recursion; (vii) predicativity and ramified
hierarchies in foundations of mathematics. Each thread is summarized
and our work positioned within it. A final subsection states the gap
this paper fills.

\section{Step-indexing in programming language semantics}

Step-indexed logical relations were introduced by Appel and McAllester
\cite{appel-mcallester2001} to give a semantic account of recursive
types in the setting of foundational proof-carrying code. The
technique gives a meaning to recursive types not via a fixed point
construction in some semantic domain, but by approximating the
intended meaning at every finite ``observation depth'' $n$. Two terms
are equivalent at index $n$ if they cannot be distinguished by any
context running for at most $n$ steps; the desired equivalence is then
the limit (intersection) of all $n$-equivalences. Appel and McAllester
show that this gives a sound model for recursive types and that the
resulting logical relation is preserved under reduction --- the
canonical \emph{step-down} property that we will recognise as the
monotonicity of step-indexed propositions.

Ahmed's PhD thesis \cite{ahmed2004} extended step-indexing to languages
with mutable references, including general references that can store
functions of any type (``higher-order store''). This is the classical
setting in which step-indexing is essential rather than merely
convenient: the standard set-theoretic recursion on type structure is
blocked by the circularity that comes from storing higher-order
functions in mutable cells. Ahmed shows how the step-indexed approach
absorbs this circularity into the index hierarchy.

Birkedal and collaborators have extended these ideas in two
directions. In one direction, \emph{topos-theoretic} reformulations
(see Section~\ref{sec:topos} below) replaced the explicit step indices
with categorical structure. In the other, the technique was scaled up
to richer programming languages and verification settings: see
Birkedal et al.\ \cite{birkedal-et-al-2011-stepindexed} on
step-indexed Kripke models for capability calculi.

The current state of the art on the programming-languages side is the
\emph{Iris} family of separation logics
\cite{jung-et-al-2018,krebbers-jung-bizjak-birkedal-dreyer-2017}.
Iris extends step-indexed separation logic with higher-order ghost
state, invariants, and a powerful set of derived rules. Its
soundness theorem is mechanized in Coq against a model built
ultimately on step-indexed propositions. The internal logic of Iris
is itself a step-indexed modal logic with a ``later'' modality
$\triangleright$ and a Löb rule; substantial libraries of Coq
tactics support reasoning within it. Other recent step-indexed
systems include VST \cite{appel2014vst}, Perceus / Koka's verified
fragment, and a number of dependent-type-theoretic frameworks
\cite{spies2021transfinite}.

Two technical observations from this literature are essential for our
work. First, the ``later'' modality has been a unifying device since
at least Appel et al.\ \cite{appel-melham-mcallester2007}, who
introduced it as the very-modal $\triangleright$ for guarded
recursion. Second, the soundness proof of every step-indexed logic
relies on a Löb-style argument --- either as an explicit lemma in the
internal logic, or implicitly through induction on the step index in
the soundness proof. Our paper takes seriously the question of which
formulation is more natural in which setting, and argues that for
higher-order recursion the internal Löb form is strictly preferable.

\section{Forcing in set theory: ramified and unramified}

Cohen introduced forcing in his 1963 papers
\cite{cohen1963-independence-i,cohen1963-independence-ii} to prove
the independence of the continuum hypothesis from \textsf{ZFC}, and
gave a book-length treatment in 1966 \cite{cohen1966}. The original
presentation is what we will call \emph{ramified forcing}: the
construction of the generic extension $M[G]$ is stratified by
ordinal levels mimicking Gödel's constructible hierarchy. A
``ramified language'' $\mathcal{L}^*$ has a type symbol for each
ordinal $\alpha$ below some sufficient bound, and a variable of type
$\alpha$ ranges only over sets at levels strictly below $\alpha$. The
forcing relation $p \Vdash_\alpha \varphi$ is defined level by level:
forcing of a level-$\alpha$ formula recurses only on forcing at levels
strictly below $\alpha$.

The two features that make this presentation ``ramified'' are
predicative stratification (a level-$\alpha$ formula only quantifies
over levels strictly below $\alpha$) and explicit level descent (the
recursive definition of forcing is structured by ordinal predecessor).
Both features are inherited from the ramified theory of types of
Russell and Whitehead \cite{russell-whitehead-PM} and from the broader
tradition of predicative analysis \cite{weyl1918,feferman1998}.

Shoenfield's 1971 lecture \cite{shoenfield1971} reformulated forcing
without the ramification: forcing conditions and $P$-names are
defined by transfinite recursion on rank, and the forcing relation is
defined by recursion on formula complexity, without explicit
stratification by syntactic level. Shoenfield's presentation is
mathematically equivalent to Cohen's --- the same generic extensions
are produced --- but it is simpler, and it has become the standard
modern presentation in textbooks \cite{kunen2011,jech2003}.
Boolean-valued models, due to Scott and Solovay
\cite{scott-solovay1967}, give yet another equivalent presentation.

What is preserved across these reformulations is the underlying
\emph{forcing relation} $p \Vdash \varphi$ together with its
fundamental properties: it is monotone (stronger conditions force
more), it is decided by some condition (every formula or its negation
is forced by some condition), and the truth of formulas in the
generic extension is captured by the forcing relation in the ground
model. What was eliminated --- the explicit ramified syntactic
hierarchy --- is exactly what we will argue is \emph{not}
eliminable in the step-indexed setting.

\section{Topos-theoretic models of step-indexing}
\label{sec:topos}

The connection between step-indexing and category-theoretic
constructions was made explicit by Birkedal, Møgelberg,
Schwinghammer, and Støvring \cite{birkedal-et-al-2012} via the
\emph{topos of trees}, the presheaf topos on the partially ordered
set $(\omega, \le)$. They observed that step-indexed propositions are
the propositions of the internal logic of this topos, and that the
``later'' modality $\triangleright$ arises naturally as a particular
endofunctor on objects of the topos. Their paper introduces
\emph{synthetic guarded domain theory}, in which all of the
step-indexed semantics is internalized: rather than indexing
manually, one works in a type theory whose intrinsic semantics is the
topos of trees.

The presheaf-topos perspective makes precise much of what the early
step-indexed literature handled informally. In particular, the
monotonicity condition on step-indexed propositions corresponds
exactly to the presheaf condition (functoriality), and the soundness
of the Löb rule is the existence of a fixed point of the
$\triangleright$ functor (the unique fixed point of $\triangleright X
= F(X)$ for a contractive $F$, or equivalently the well-foundedness
of $(\omega, <)$ as the index domain). Subsequent work has refined
the categorical setting and used it to model more sophisticated
languages \cite{bizjak-grathwohl-clouston-mogelberg-birkedal-2016}.

The presheaf-topos perspective is implicitly a forcing-style semantics
in the model-theoretic sense: a presheaf topos on $(P, \le)$ is
exactly the topos of sheaves with respect to the trivial Grothendieck
topology on $P$, and these are exactly the Kripke models of
intuitionistic logic with worlds in $P$. The topos's truth value
object plays the role of a Heyting algebra-valued ``forcing'' truth
value. To our knowledge, this connection has not been emphasized in
the topos-of-trees literature, and the word ``forcing'' has not been
prominent in it. Our paper makes this terminological and conceptual
connection explicit.

\section{Forcing translations in type theory}
\label{sec:jts}

A parallel literature, originating in the constructive type theory
community, develops syntactic \emph{forcing translations} of type
theories. The seminal work here is Jaber, Tabareau, and Sozeau
\cite{jaber-tabareau-sozeau-2012}, who give a forcing translation of
the Calculus of Constructions: given a forcing notion (a poset
$P$), they exhibit a syntactic translation that takes a CIC term
$M : T$ to a forced term $M^P : T^P$, where the type translation
$T \mapsto T^P$ corresponds to ``the forced version of $T$.'' For
propositions, $T^P$ is a downward-closed predicate on $P$ ---
exactly our \texttt{iProp} of Section~\ref{sec:prop-layer}. The
Jaber et al.\ translation is computational: it gives a definitional
account of forcing within type theory itself, not just a semantic
interpretation.

This line was extended in
\cite{jaber-et-al-2016} (the ``definitional side of the forcing''),
which gives a refined translation that preserves more of the
computational structure of the original type theory. Coquand and
collaborators have related work on constructive presentations of
forcing in type-theoretic settings \cite{coquand-jaber2010}.

The Jaber et al.\ translations are the type-theoretic version of
ramified forcing for type theory: a parametric construction of a
``forced universe'' that supports the same operations as the
unforced one. Our Phase 4 abstract framework recovers the
propositional fragment of this translation by direct construction in
Rocq. The Phase 6 file makes the embedding precise via a constant
embedding $\mathsf{Prop} \hookrightarrow \mathtt{iProp}$ that is
a Heyting algebra homomorphism.

What our work does not do --- and what the Jaber et al.\ line does ---
is give a translation of the \emph{full} type theory, including
dependent types. A natural future direction is to formalize the
restriction of the Jaber et al.\ translation to the propositional
fragment in Rocq and check (mechanically) that it produces our
\texttt{iProp} up to definitional equality.

\section{The modal logic of provability and well-foundedness}
\label{sec:gl}

The Gödel--Löb provability logic \textsf{GL} is the modal logic
characterized by the axiom $\Box(\Box \varphi \to \varphi) \to
\Box \varphi$ (the ``Löb axiom''), together with the standard
$\mathsf{K}$ axioms. Solovay's arithmetical completeness theorem
\cite{solovay1976} establishes that \textsf{GL} is sound and complete
for arithmetical interpretations of $\Box$ as ``it is provable in
Peano Arithmetic that\ldots''. From the model-theoretic side,
\textsf{GL} is sound and complete for the class of transitive
Kripke frames whose accessibility relation $R$ is
\emph{converse well-founded} (no infinite ascending chain
$w_0 R w_1 R w_2 R \cdots$); see Smoryński \cite{smorynski1985} and
Boolos \cite{boolos1993} for textbook treatments.

The model theory of \textsf{GL} is directly relevant to our work.
The ``later'' modality $\triangleright$ of step-indexed logic is
formally exactly a \textsf{GL}-style provability $\Box$ once we
identify the strict refinement order $\prec$ on conditions with the
converse of the \textsf{GL} accessibility relation. Löb's rule of
step-indexed logic is the Löb axiom of \textsf{GL}. The fact that
$\triangleright$ requires well-foundedness for Löb's rule to hold is
the dual of \textsf{GL}'s requirement that the accessibility relation
be converse well-founded. Our Phase 5 result --- that Löb fails over
$\mathbb{Z}$ --- is therefore a special case of (the converse of)
the soundness of \textsf{GL} for converse-well-founded frames.

What our paper adds to this picture is the observation that the
familiar PL-style step-indexed logic, including its mechanizations in
Iris and elsewhere, is in fact a propositional fragment of \textsf{GL}
indexed over a particular kind of \textsf{GL}-frame (one with
well-founded strict refinement). This connection is folklore in modal
logic but, as far as we are aware, has not been made explicit or
mechanized for the step-indexed Hoare logic setting.

\section{Constructive modal logic and guarded recursion}

A separate line of work, with substantial overlap with synthetic
guarded domain theory, develops constructive modal logics for
guarded recursion. Atkey and McBride \cite{atkey-mcbride2013} use the
$\triangleright$ modality to characterize \emph{productive
coprogramming}: a corecursive function is productive iff each recursive
call sits under a $\triangleright$. Clouston, Bizjak, Grathwohl, and
Birkedal \cite{clouston-et-al-2015,clouston-et-al-2017} develop
guarded $\lambda$-calculi with $\triangleright$ as a primitive type
former, supporting recursive types via a guarded fixed-point operator.
Møgelberg and others \cite{mogelberg2014,bizjak-mogelberg-2018}
extend these systems to dependent types and provide categorical
semantics.

This literature is closely allied with synthetic guarded domain
theory: both treat $\triangleright$ as a built-in modal operator and
exploit its compatibility with recursion. The work is structured
mostly around \emph{computational} and \emph{type-theoretic} questions
(``how do we type recursion productively?''), rather than
\emph{semantic} questions (``what does the step-indexed semantics
look like as a forcing construction?''). Our paper is complementary:
we work on the semantic side, with the forcing reading as our
organizing principle.

\section{Predicativity and ramified hierarchies}

The word ``ramified'' belongs to a tradition older than Cohen's
forcing. Russell and Whitehead \cite{russell-whitehead-PM} introduced
the ramified theory of types in \emph{Principia Mathematica} to
resolve the set-theoretic paradoxes via a stratification of
propositions by ``order,'' where a proposition of order $n$ may
quantify only over propositions of order $< n$. The same idea
animates Weyl's predicative analysis \cite{weyl1918}, in which the
real numbers are constructed within an arithmetical hierarchy without
appeal to impredicative set-theoretic principles. Feferman's program
of \emph{predicativity given the natural numbers} \cite{feferman1998}
develops this into a comprehensive foundation for substantial parts of
classical mathematics.

The thread connecting ramified type theory to ramified forcing is
that both stratify quantification by a well-founded hierarchy, and
both treat the hierarchy itself as part of the syntactic apparatus
rather than as a meta-level convenience. Shoenfield's
unramification of forcing showed that for forcing the stratification
can be absorbed into a single recursion on rank; what we argue, by
mechanizing step-indexed Hoare logic in higher-order settings, is
that for step-indexing the analogous unramification \emph{cannot}
be performed. The well-founded forcing notion is doing real work,
not just structuring the proof.

This connection also clarifies why the Iris-style internal logic
``feels predicative'': step-indexed propositions are stratified by a
well-founded hierarchy, and any recursive definition that does not
respect the hierarchy --- by mediating recursion through a
$\triangleright$ --- is rejected. The system is not just predicative in
spirit; it is precisely a ramified predicative system in the technical
sense.

\section{Hoare logic and program verification}

We omit a survey of classical Hoare logic and separation logic,
referring the reader to standard sources
\cite{hoare1969,ohearn-reynolds-yang-2001,brookes2007}. Two points
are worth noting. First, our IMP development in Phase 2 uses the
plain (non-separation) version of Hoare logic, since separation
machinery is not needed for the conceptual point we want to make
about the \texttt{while} rule. Second, Cook's relative completeness
theorem \cite{cook1978} for Hoare logic shows that the classical
proof system is as complete as can be hoped given any reasonable
specification language; our results are orthogonal to this --- we
study how the soundness side is best organized, not the proof system.

\section{The gap this paper fills}

The summaries above suggest that almost every ingredient of our work
is, individually, well-known to specialists in some adjacent
community. The topos-of-trees perspective on step-indexing is
familiar in synthetic guarded domain theory; the \textsf{GL} modal
logic of provability is familiar to modal logicians; the
ramified-forcing presentation of Cohen is familiar to set theorists;
the JTS forcing translation is familiar to the type theory
community; and the predicative reading of ramification is familiar
to foundations researchers.

What does \emph{not} appear in the literature, to our knowledge, is
a single mechanized account of step-indexed Hoare logic that:

\begin{enumerate}

\item \textbf{Frames step-indexing explicitly as ramified forcing}
(Section~\ref{sec:thesis}), making the connection to Cohen's
original presentation precise rather than implicit in the
topos-theoretic literature.

\item \textbf{Demonstrates that internal Löb is the load-bearing
soundness tool for higher-order recursion} by mechanizing both an
IMP-style (Phase 2) and a $\lambda$-calculus-with-\texttt{fix}
(Phase 3) Hoare logic and exhibiting both the meta-level and
internal forms of the \texttt{while} rule's soundness proof.

\item \textbf{Generalizes the framework to abstract forcing notions}
parameterically (Phase 4), exhibiting at least one second
instance (the pointwise product $\omega \times \omega$, suited to
parallel components) and recovering the standard $(\omega, \le)$
instance.

\item \textbf{Proves the well-foundedness of the forcing notion is
essential} (Phase 5) by an explicit counterexample over $\mathbb{Z}$
where the iProp constructions go through but Löb's rule fails.
This is a special case of the \textsf{GL} model theory, but the
specific statement for step-indexed Hoare logic does not appear in
the PL literature we surveyed.

\item \textbf{Verifies the bridge to the JTS forcing translation}
(Phase 6) by formalizing the constant embedding
$\mathsf{Prop} \hookrightarrow \mathtt{iProp}$ as a Heyting algebra
homomorphism and identifying $\triangleright$ as the
step-indexed-specific addition.

\end{enumerate}

\noindent
The contribution is therefore not the introduction of any single new
construction but the \emph{integration} of several existing threads
into one mechanized account, plus the small new technical result of
Phase 5 and the explicit comparison of Phase 2c and Phase 3c that
diagnoses where internal Löb earns its keep.

\printbibliography

\end{document}
